% Noether Paper 40 controlled-generic zh-Hant producer draft.
% Translation source: exact current-German marker interval, SHA-256
% 7965805D3A75C3354C85BC7A3E4725F07BF869A8833FC19D74E32BE369427937.
% Inherited Simplified-Chinese interval is a drafting witness only, SHA-256
% 3DAD18CAB878BDFA62ED4FCC634E21AF92AF22BC8E11DF8A36888088D0A608AB.
% Status: controlled-generic script draft; independent check pending.
% Compilation is mechanical production only. No source, semantic, formula,
% terminology, visual, regional, publication, or certification check is claimed.
% Controlled generic Traditional script only; not zh-Hant-TW/HK/MO prose.
% Hans wording remains the lexical base; independent Hant checking is pending.
\documentclass[11pt]{article}
\usepackage[a4paper,margin=2.4cm]{geometry}
\usepackage{fontspec}
\usepackage{xeCJK}
\usepackage{amsmath,amssymb,mathtools,array,mathrsfs}
\setmainfont{Noto Serif}
\setCJKmainfont{Microsoft JhengHei}
\setCJKsansfont{Microsoft JhengHei}
\XeTeXlinebreaklocale "zh"
\XeTeXlinebreakskip=0pt plus 1pt
\setlength{\parindent}{2em}
\setlength{\parskip}{0.35em}
\setlength{\emergencystretch}{8em}
\allowdisplaybreaks
\raggedbottom
\makeatletter
\providecommand{\srcfn}[2]{%
  \begingroup
  \renewcommand{\thefootnote}{#1}%
  \footnote{#2}%
  \addtocounter{footnote}{-1}%
  \endgroup}
\providecommand{\srcfnmark}[1]{%
  \begingroup
  \renewcommand{\thefootnote}{#1}%
  \footnotemark%
  \addtocounter{footnote}{-1}%
  \endgroup}
\providecommand{\srcfntext}[2]{%
  \begingroup
  \renewcommand{\thefootnote}{#1}%
  \footnotetext{#2}%
  \endgroup}
\makeatother
\providecommand{\srcspaced}[1]{#1}
\providecommand{\glqq}{“}
\providecommand{\grqq}{”}
\providecommand{\frR}{\mathfrak R}
\providecommand{\frS}{\mathfrak S}
\providecommand{\frT}{\mathfrak T}
\providecommand{\frM}{\mathfrak M}
\providecommand{\frN}{\mathfrak N}
\providecommand{\frA}{\mathfrak A}
\providecommand{\frB}{\mathfrak B}
\providecommand{\frC}{\mathfrak C}
\providecommand{\frG}{\mathfrak G}
\providecommand{\frL}{\mathfrak L}
\providecommand{\frako}{\mathfrak o}
\providecommand{\mA}{\mathfrak A}
\providecommand{\mbarA}{\overline{\mathfrak A}}
\providecommand{\Gcal}{\mathfrak G}
\providecommand{\Sp}{\operatorname{Sp}}
\begin{document}
\setcounter{footnote}{0}

\section*{40. 非交換代數}

\begin{center}
\emph{Math. Zs. 37 (1933), 514--541}
\end{center}

\begin{center}
\textbf{非交換代數。}\\[0.5em]
作者\\[0.25em]
哥廷根的 Emmy Noether。
\end{center}

眾所周知，交換代數的主要定理都包含在伽羅瓦理論之中；在此之前則有析出域與分裂域的理論，也就是關於這樣一些域的理論：它們足以使一個給定多項式分出一個一次因子，或分別使它完全分解為一次因子。

本文把代數的相應部分發展到非交換情形，尤其發展到超復情形。而且我在這裡原則上採用非交換方法，即在非交換除環中的表示。最後我將說明，上述交換代數定理如何能夠完全平行地通過在交換域中的表示來加以證明。

這裡所依據的表示論之前有一段簡短的自同態理論（§ 1）；它是對以表示模理論為基礎的表示論的進一步發展\footnote{E. Noether, Hyperkomplexe Größen und Darstellungstheorie, Math. Zeitschr. 30 (1929), S. 641--692, 下文簡稱“表示論”。另參見 van der Waerden, Moderne Algebra II 中的轉述。}。尤其是，除直接表示之外，我還考察反向表示——其中一種可以歸結為另一種——；現在，反向表示由反向表示模產生。其優點在於，反向表示模——也就是一個雙模——還可以理解為擴張環上的單側模（§ 2）。這樣一來，非交換除環中的表示便歸結為擴張環中的理想論；不可約表示類對應於擴張環的不可約理想類，這與超復系統在其交換系數域中的表示對系統本身所成立的事實完全類似（§ 3 和 4）。

由此，藉助如下說明（§ 5）便得到非交換除環上矩陣環的結構定理：\emph{每一個子環都給出這個除環上的一個表示，或分別給出由其反向同構除環所作的反向表示}。這個反向同構除環構成交換情形中極小析出域與分裂域的第一個類似物，因為它給出所有一次反向表示；而當它關於中心的秩有限時——此前並未作這種假設——還給出分解為秩為 1 的直和項的完全分裂。由此得到除環的伽羅瓦理論（§ 6），同時也表達瞭如下事實（§ 7）：反向同構的除代數（關於中心具有有限秩的除環）在 R. Brauer 的代數類群中生成互逆的類。

伽羅瓦理論還有第二種證明，它更一般地適用於簡單系統（並在正文中先行給出）；它幾乎直接來自一個關於兩兩可交換子環的定理，而這個定理本身又幾乎直接銜接於前置的自同態考察。

截至這裡，全部考察都是純粹非交換的。不過，對\emph{交換}分裂域的提問也用非交換方法處理：依據 § 5 中預先給出的說明，由除代數來表示分裂域（§ 7）。最後是上述向交換情形的轉移（§ 8），以及任意系統的分裂域理論（§ 9）；與此同時也就給出了它同交換情形中通常的表示論之間的聯繫。

R. Brauer 正是把分裂域理論建立在這種交換情形中的表示論以及“無理”因子系之上\footnote{參見我們的合著札記：Über minimale Zerfällungskörper irreduzibler Darstellungen. Sitz. Ber. d. Preuß. Ak. d. Wiss. 1927, S. 221--228。（第 222 頁的註釋含有一份“說明報告”。主要定理是彼此獨立而且幾乎同時產生的。）另參見 R. Brauer, Über Systeme hyperkomplexer Zahlen, § 3. Math. Zeitschr. 30 (1929), S. 79--107。A. A. Albert 後來也獨立地重新發現了這些定理。}。由於這種交換的證明，他——以及後來的 Albert——必須假定中心是完美域；從非交換證明來看，這是一個不必要的限制。在同一限制下，同樣藉助交換情形中的表示論，並且是在得知我的非交換除環伽羅瓦理論之後，R. Brauer 和 K. Shoda 又進一步發展了這一理論：R. Brauer 給出了上述關於兩兩可交換子環的定理，K. Shoda 則獨立地給出了也適用於半單系統的完整理論\footnote{R. Brauer, Über die algebraische Struktur von Schiefkörpern, § 2. Journ. f. Math. 166 (1932), S. 241--252. K. Shoda, Über die galoissche Theorie der halbeinfachen hyperkomplexen Systeme. Math. Ann. 107 (1932), S. 252--258。J. Levitzki 也曾發展出這些結果。}。

最後還應提到，對於除環的情形，van der Waerden II 中有一個簡短的論述（§ 128），它接續於我 1928 年夏季學期的課程；我在那裡第一次講授了這些問題。van der Waerden 的論述比這門課的講法簡化了許多地方，其中一部分我在第二門課程（1929/30 年冬季學期）\footnote{第一門課程由 G. Köthe 整理，第二門由 M. Deuring 整理。}中採用，另一部分則直到本文才採用並繼續發展。特別是，把伽羅瓦理論的超復證明方法轉移到交換情形（§ 8），最早來自第二門課程；在那裡，非交換證明（§ 6, 3.）比第一門課程中的證明大為簡化。關於兩兩可交換子環的定理（§ 5），以及以此為基礎的非交換伽羅瓦理論的第二種證明方法（§ 6, 1. 和 2.），是在第二門課程之後、得知 R. Brauer（注 3））和 van der Waerden § 128 的工作後才加入的；不過 § 5, 3. 中所需的有限性假設比那些工作弱得多。

第一門課程中的某些推理方法——對差分理想取交的方法——雖然較為複雜，卻有一個優點，即可以轉移到無限秩系統和整系統上\footnote{這一方法實質上見於 G. Köthe, Schiefkörper unendlichen Ranges über dem Zentrum § 5, Math. Ann. 105 (1931), S. 15--39。另見注 10）。這種取交方法現在已被一個幾乎平凡的事實（§ 4, 1.）所取代——這個事實最早由 van der Waerden 注意到——，即雙側理想屬於不變模。這樣一切都能通過更簡單的直和分解來證明；但這一方法在無限秩和整情形中失效。}。對於交換的整系統，由此可以得到理想微分與不同之間的聯繫\footnote{參見布拉格的報告，Jahresber. d. Deutsch. Math. Ver. 39 (1930), S. 17（斜向分頁）。}，我將在別處更詳細地討論這一點。

分裂域理論的主要應用是在交叉積及其因子系的理論中；後者又構成數論應用的基礎。本文不再討論這一方面\footnote{交叉積理論是在第二門課程中發展出來的；經過少量修改，使之不必預設本文給出的理論，該理論收入 H. Hasse: Theory of cyclic algebras over an algebraic numberfield, Kap. II. Transactions of the Amer. Math. Soc. 34 (1932), S. 171--214。M. Deuring 將在“Ergebnisse der Mathematik”的一篇報告中給出一種更貼近課程講稿的論述。}。

\begin{center}
\textbf{§ 1.}\\[0.4em]
\textbf{自同態、模與雙模。}
\end{center}

眾所周知，以表示模為基礎的表示論（另見 § 2）依賴於帶算子或不帶算子的阿貝爾群之自同態環理論。為避免重複，下面把其中隱含採用的推理——映射與運算法則之間的關係——表述成幾個簡單命題。

\noindent\textbf{1. 乘法映射。結合律。} 首先設 \(\frG\) 是一個不帶算子的群，\(\frA\) 是它的絕對自同態系統，即從 \(\frG\) 到自身的所有同態所組成的系統。\(\frA\) 對乘法封閉，因為乘積 \( \sigma\tau \) 由
\[
        g(\sigma\tau)=(g\sigma)\tau
        \quad \text{其中 } g \text{ 屬於 } \frG \text{，且 } \sigma,\tau \text{ 屬於 } \frA                         \tag{1}
\]
定義，並且顯然滿足結合律。

眾所周知，如果給定一個由符號 \(\Theta,H,\ldots\) 組成的集合 \(\frB\)，使得運算 \(g\Theta,gH,\ldots\) 給出 \(\frG\) 中唯一確定的元素，併產生 \(\frG\) 的自同態（到自身的同態）——\((g\cdot h)\Theta=g\Theta\cdot h\Theta\)——，那麼 \(\frG\) 就成為一個帶算子的群；因此，從 \(\frB\) 到絕對自同態系統 \(\frA\) 的一個子集 \(\overline{\frB}\) 存在唯一的映射（一般並非雙射）。上述推理在這裡表述如下：

\emph{若算子域 \(\frB\) 對乘法封閉，則從 \(\frB\) 到 \(\overline{\frB}\) 的唯一映射是乘法同態，當且僅當滿足與 (1) 相應的結合關係}
\[
        g(\Theta H)=(g\Theta)H
        \quad \text{其中 } g \text{ 屬於 } \frG \text{，且 } \Theta,H \text{ 屬於 } \frB                         \tag{1a}
\]
\emph{。相應地，在處理左算子時，所定義的是反向同態。}

這是因為，從 \(\frB\) 到 \(\overline{\frB}\) 的映射由對 \(\frG\) 中所有 \(g\) 成立的 \(g\Theta=g\sigma\) 給出，因而也由
\[
        (g\Theta)H=(g\sigma)\tau=g(\sigma\tau)
\]
給出，所以 (1a) 成為乘法同態的充要條件。左算子之所以產生反向同態，是因為寫在左邊的自同態 \(\sigma^\ast g,\tau^\ast\sigma^\ast g\) 必須從右向左讀：\(\tau^\ast\sigma^\ast g=g\sigma\tau\)，而算子則總是從左向右讀。

\noindent\textbf{2. 算子同態映射。} 若 \(\frG\) 是帶算子的群，則眾所周知，算子同態性由
\[
        \begin{array}{r@{\quad}c@{\quad}c@{\quad}c@{\qquad}c@{\qquad}r@{\quad}c@{\quad}c@{\quad}c}
        (2) & (g\Theta)\sigma &=& (g\sigma)\Theta
        & \text{或} &
        (2*) & (\Theta g)\sigma &=& \Theta(g\sigma),
        \end{array}
\]
也就是分別由交換相容性或貫通結合律來定義。對兩個不同的算子域採用 1. 中的推理——映射到自同態系統——便得到：

\emph{若給定兩個算子域 \(\frB\) 與 \(\frC\)，其元素分別為 \(\Theta,H,\ldots\) 與 \(\bar\Theta,\bar H,\ldots\)，則 \(\frC\) 的元素產生 \(\frB\)-自同態、同時 \(\frB\) 的元素產生 \(\frC\)-自同態，當且僅當這些域同 \(\frG\) 交換相容，或分別滿足貫通結合律：}
\[
        \begin{array}{r@{\quad}c@{\quad}c@{\quad}c@{\qquad}r@{\quad}c@{\quad}c@{\quad}c}
        (2a) & (g\Theta)\bar H &=& (g\bar H)\Theta, &
        (2a*) & (\Theta g)\bar H &=& \Theta(g\bar H).
        \end{array}
\]

\noindent\textbf{3. 環上的模與雙模。} 眾所周知，環 \(\frR\) 上的一個\emph{右模} \(\frM\) 定義為一個加法阿貝爾群，以 \(\frR\) 中的元素作為右算子；除了結合關係 (1a) 與定義算子的關係 \((g+h)\Theta=g\Theta+h\Theta\) 之外，它還滿足分配關係
\[
        \begin{array}{r@{\quad}c@{\quad}c@{\quad}c}
        (3a) & g(\Theta+H) &=& g\Theta+gH
        \end{array}
\]
；左模也相應定義。

\emph{雙模}須分成兩種：環 \(\frR\) 與 \(\frS\) 上的右模稱為雙模，如果除了分別對 \(\frR\) 與 \(\frS\) 成立的運算規則 (1a, 3a) 之外，\(\frR\) 與 \(\frS\) 還依據 (2a) 同 \(\frM\) 交換相容；\(\frR\)-左、\(\frS\)-右模稱為雙模，如果在其餘運算規則之外還滿足 (2a*)，即貫通結合律。

這些運算規則的意義來自如下事實：在\emph{阿貝爾}群的情形，絕對自同態系統構成一個環，即絕對自同態環\footnote{對非阿貝爾群，所涉及的是一個“廣義”環；參見 H. Fitting 一篇將發表於 Math. Annalen 的論文。[Math. Annalen 107, S. 514--542.] }。具體地，映射推理給出：

\emph{若一個（以加法記號書寫的）阿貝爾群 \(\frM\) 的算子域 \(\frR\) 是一個環，則從 \(\frR\) 到絕對自同態環的一個子集 \(\overline{\frR}\) 的唯一映射是環同態，當且僅當 \(\frM\) 是 \(\frR\) 上的右模——即滿足 (1a) 與 (3a)——；對左模則是反向環同態。}

\emph{若 \(\frM\) 是環 \(\frR\) 與 \(\frS\) 上的右模，則 \(\frM\) 成為雙模——即滿足 (2a)——，當且僅當 \(\frR\) 可以環同態地映到 \(\frS\)-自同態環的一個子環上，同時 \(\frS\) 可以環同態地映到 \(\frR\)-自同態環的一個子環上。若 \(\frM\) 是 \(\frS\) 上的右模、\(\frR\) 上的左模，則雙模條件——即 (2a*)——意味着 \(\frS\) 同態地映到 \(\frR\)-自同態環的一個子環上，而 \(\frR\) 反向同態地映到 \(\frS\)-自同態環的一個子環上。}

最後一項陳述使下列已知推論得到更精確的表述：

\emph{1) 若 \(\frR\) 是有單位元的環，則作為 \(\frR\)-左模，\(\frR\) 與其自同態環直接同構——而不僅僅同態；作為 \(\frR\)-右模，則與其自同態環反向同構。} 這是因為貫通結合律 (2a*) 成立，所以同態陳述成立。單位元的存在使每一個自同態都由 \(\frR\) 中的對應 \(e\mapsto a\) 給出；因而這裡得到的是同構，並且 \(\frR\) 窮盡這個 \(\frR\)-自同態環。

\emph{2) 若 \(\frR\) 是一般非交換除環 \(A\) 上的全矩陣環，}
\[
        \frR=\sum c_{ik}A=\sum A c_{ik},
\]
\emph{則 \(A\) 與單右理想的自同態除環反向同構，而與單左理想的自同態除環直接同構。} 這是因為每一個單右理想都與一個 \(A\)-左、\(\frR\)-右模算子同構，並且 \(A\) 窮盡這些 \(\frR\)-自同態；單左理想的情形相應成立。

\noindent\textbf{4. 從右雙模過渡到乘積環上的單側模。} 右雙模的重要性實質上正是以這種過渡為基礎。

\noindent\textbf{定理：} \emph{若 \(\frM\) 是環 \(\frT\) 上的右模，而 \(\frT\) 包含兩個元素兩兩可交換的子環 \(\frR\) 與 \(\frS\)，則 \(\frM\) 也可以理解為 \(\frR\) 與 \(\frS\) 上的雙模。反過來，若給定 \(\frR\) 與 \(\frS\) 上的一個雙模，並且存在乘積環 \(\frT\)，在其中 \(\frR\) 與 \(\frS\) 的元素兩兩可交換，則 \(\frM\) 也可以理解為 \(\frT\) 上的右模，使得 \(\frT\) 的運算是給定的 \(\frR\) 與 \(\frS\) 運算的延拓。}

第一項陳述來自如下事實：按照定義，\(\frT\) 可以同態地映到絕對自同態環的一個子環 \(\overline{\frT}\) 上，其中 \(\frR,\frS\) 分別變為元素兩兩可交換的子環 \(\overline{\frR},\overline{\frS}\)。而這恰好就是刻畫雙模的關係 (2a)。

反過來，若給定 \(\frM\) 為 \(\frR\)、\(\frS\)-雙模（右模），則 \(\frR,\frS\) 仍可同態地映到絕對自同態環中元素兩兩可交換的子環 \(\overline{\frR},\overline{\frS}\) 上。在絕對自同態環中，對 \(\overline{\frR},\overline{\frS}\) 存在乘積環\footnote{“乘積環”在這裡只是指由兩個具有公共上環的環所生成的環；這個乘積不必是直積。——給定交時，乘積環未必總是存在，下例說明這一點：設 \(\frako\) 為整數環，並令 \(\frR=\frako[x]\)、\(\frS=\frako[y]\)，其中 \(2x-1=0\)、\(2y-1=0\)；若未規定 \(x\) 與 \(y\) 之間的運算關係，則 \(\frako\) 是 \(\frR\) 與 \(\frS\) 的交。但若 \(\frR\) 與 \(\frS\) 位於一個公共上環中，則有：\((2x-1)y-x(2y-1)=0\)，從而 \(x=y\)。因此不存在以 \(\frako\) 為交的乘積環。} \(\overline{\frT}\)；由於元素兩兩可交換，它由所有形如
\[
        \sum \bar r_i\bar s_i+\bar r+\bar s
\]
的元素組成（若 \(\overline{\frR},\overline{\frS}\) 有單位元，則附加項 \(\bar r,\bar s\) 可以省去）。如果還存在一個乘積環 \(\frT\)，其中 \(\frR\) 與 \(\frS\) 的元素兩兩可交換，那麼它相應地由所有形如 \(\sum r_i s_i+r+s\) 的元素組成。因此，同態
\[
        \frR\to\overline{\frR},\qquad \frS\to\overline{\frS}
\]
可以通過對應
\[
        \sum r_i s_i+r+s
        \longmapsto
        \sum \bar r_i\bar s_i+\bar r+\bar s
\]
擴充為一個同態 \(\frT\to\overline{\frT}\)；所以運算
\[
        m\Bigl(\sum r_i s_i+r+s\Bigr)
        =\sum (m r_i)s_i+m r+m s
\]
是唯一確定的，並按照 §1, 3. 定義一個包含給定 \(\frR\)-、\(\frS\)-模的 \(\frT\)-模。

\begin{center}
\textbf{§ 2.}\\[0.4em]
\textbf{反向表示與直接表示。}
\end{center}

\noindent\textbf{1. 線性形式模。} 從 §1 所發展的理論過渡到表示論，是以把其中的模專門化為線性形式模並考察其自同態環為基礎的。

一個 \(\frS\)-右模 \(\frM\)——其中 \(\frS\) 是有單位元的環——稱為 \(\frS\) 中的線性形式模，如果 \(\frM\) 是 \(n\) 個單生成元 \(\frS\)-模的直和：\(\frM=m_1\frS+\cdots+m_n\frS\)，並且每個 \(m_i\frS\) 都同 \(\frS\) 算子同構，也就是從 \(m_i s=0\) 總能推出 \(s=0\)。由此可知，\(\frS\) 不僅可以同態地，而且可以同構地映到絕對自同態環的一個子環上。

\noindent\textbf{定理 1。} \emph{若 \(\frM\) 是 \(\frS\) 中的線性形式模，則 \(\frM\) 的 \(\frS\)-自同態環 \(\mA\) 與 \(\frS\) 中所有 \(n\) 階矩陣組成的環 \(\mbarA\) 反向同構——而不僅僅同態。}

\(\frM\) 的任意一個 \(\frS\)-自同態 \(\alpha\) 都由對應 \(m_i\longmapsto \overline m_i=m_i\alpha\) 完全描述；因為由定義可得
\[
        \sum m_i s_i\longmapsto
        \sum \overline m_i s_i
        =\sum (m_i\alpha)s_i
        =\sum (m_i s_i)\alpha.
\]
若現在用矩陣記法寫成
\[
        (m_1\alpha,\ldots,m_n\alpha)=(m_1,\ldots,m_n)A,
\]
那麼，當 \(\alpha\) 遍歷所有自同態時，對應 \(\alpha\mapsto A\) 給出 \(\mA\) 與 \(\mbarA\) 之間的一一對應。因為已經假定 \(\frM\) 為線性形式模，所以 \(A\) 由 \(\alpha\) 唯一確定；並且 \(\frS\) 中每一個 \(n\) 階矩陣 \(A\) 都產生一個自同態。此外還有：
\[
        \alpha+\beta\mapsto A+B,\qquad
        \alpha\beta\mapsto BA.
\]
後一個對應來自
\[
        (m_1\alpha\beta,\ldots,m_n\alpha\beta)
        =(m_1,\ldots,m_n)A\beta
        =(m_1,\ldots,m_n)\beta A
        \quad[\text{因為 }m s\beta=m\beta s]
        =(m_1,\ldots,m_n)BA.
\]

\noindent\textbf{2. 表示與表示模。} 如果把環 \(\frR\) 在 \(\frS\) 中的 \(n\) 階反向表示或直接表示，分別理解為從 \(\frR\) 到 \(\frS\) 中所有 \(n\) 階矩陣所成環的某個子環的反向環同態或直接環同態，那麼定理 1 也可表述為：

\noindent\textbf{定理 1'。} \emph{\(\frS\) 中一個 \(n\) 階線性形式模的自同態環 \(\mA\)，可以由 \(\frS\) 中所有 \(n\) 階矩陣組成的全矩陣環 \(\mbarA\) 給出一個同構的（忠實的）反向表示。}

由此，直接表示所通常具有的定理也同樣可用於反向表示。

\noindent\textbf{定義。} \(\frS\) 中的一個線性形式模 \(\frM\) 稱為 \(\frR\) 在 \(\frS\) 中的反向表示模，如果 \(\frM\) 是 \(\frR\) 與 \(\frS\) 上的雙模，同時又是 \(\frR\)-右模和 \(\frS\)-右模；它稱為直接表示模，如果 \(\frM\) 是雙模，同時又是 \(\frR\)-左、\(\frS\)-右模。

\noindent\textbf{定理 2。} \emph{每一個反向表示模或直接表示模，都產生 \(\frR\) 在 \(\frS\) 中的一類等價的反向表示或直接表示；而所有反向表示或直接表示都是以這種方式產生的。}

首先設 \(\frM\) 是反向表示模。按照 §1, 3.，\(\frR\) 可以直接同態地映到 \(\frM\) 的 \(\frS\)-自同態環的一個子環 \(\overline{\frR}\) 上；按照 §2, 定理 1'，\(\overline{\frR}\) 具有一個反向的（忠實的）表示，因而它也成為 \(\frR\) 的反向表示。反過來，若給定一個反向表示 \(\frR\to\frR^*\)，把它與從 \(\frR^*\) 到 \(\frS\)-自同態環的一個子環 \(\overline{\frR}\) 的反向同構複合，便得到從 \(\frR\) 到 \(\overline{\frR}\) 上的直接同態；所以 \(\frM\) 成為一個 \(\frR\)-模，而且成為雙模，也就是按照 §1, 3. 的表示模。換用別的 \(\frS\)-基便得到等價表示類。

若討論的是直接表示模，則 \(\frR\) 被反向同態地映到 \(\overline{\frR}\) 上。與 \(\overline{\frR}\) 的反向表示複合後，便得到 \(\frR\) 的一個直接表示。反過來，從一個直接表示可得到到 \(\overline{\frR}\) 上的反向同態映射；於是按照 §1, 3.，\(\frM\) 成為一個直接表示模。

\noindent\textbf{說明。} 兩種表示當然可以互相歸結：\(\frR\) 的一個直接表示就是與 \(\frR\) 反向的環的一個反向表示，反之亦然。另一種歸結方式是從 \(\frS\) 過渡到一個反向環，並以轉置矩陣代替原矩陣。這對應於從一個 \(\frR\)-左、\(\frS\)-右模過渡到一個 \(\frR\)-左、\(\frS\)-左模。

\noindent\textbf{3. 從反向表示模過渡到擴張環上的模。} 反向表示模的優點首先基於 §1, 4. 所允許的向擴張環的過渡。對於這裡的特殊情形——\(\frM\) 是 \(\frS\) 上的線性形式模——，也就是表示模的如下過渡定理：

\emph{若 \(\frM\) 是環 \(\frT\) 上的右模，而 \(\frT\) 包含兩個元素兩兩可交換的子環 \(\frR\) 與 \(\frS\)，並且 \(\frM\) 成為 \(\frS\) 上的線性形式模，則 \(\frM\) 也可以理解為從 \(\frR\) 到 \(\frS\) 的反向表示模。}

\emph{反過來，若給定一個從 \(\frR\) 到 \(\frS\) 的反向表示模，並且存在一個乘積環 \(\frT\)，其中 \(\frR\) 與 \(\frS\) 的元素兩兩可交換，則 \(\frM\) 也可以理解為一個 \(\frT\)-模，使得 \(\frT\) 的運算是給定的 \(\frR\)-、\(\frS\)-運算的延拓。}

表示論中實質上使用的是

\noindent\textbf{表示模過渡定理的推論。}\\*
若 \(\frM\) 是 \(\frR\) 在 \(\frS\) 中的反向表示模，同時又是以 \(\frT\) 為乘積環的 \(\frT\)-模，則 \(\frM\) 到某個模 \(\frN\) 的 \(\frR\)-、\(\frS\)-同構等價於 \(\frT\)-同構。\emph{因而，每一類同構的 \(\frT\)-模都對應於 \(\frR\) 在 \(\frS\) 中的一個反向表示類，反之亦然。}

這些定理同 §1, 2. 和 3. 相結合，給出下面所需的基本聯繫，即與某個表示可交換的矩陣同 \(\frR\)-、\(\frS\)-自同態之間的聯繫。

\noindent\textbf{關於可交換矩陣的定理。}\\*
\emph{若 \(\frR\to\frR^*\) 是 \(\frR\) 在 \(\frS\) 中的一個 \(n\) 階反向表示，並且 \(\frB^*\) 表示所有與 \(\frR^*\) 的元素兩兩可交換的矩陣（\(\frS\) 中的 \(n\) 階矩陣）所成的環，則 \(\frB^*\) 給出生成該表示的表示模之 \(\frR\)-、\(\frS\)-自同態的一個反向同構表示——也就是那些同時還是 \(\frR\)-自同態的 \(\frS\)-自同態之反向同構表示——；若乘積環 \(\frT\) 存在，也就是 \(\frT\)-自同態的反向同構表示。}

設 \(\frB\) 是 \(\frM\) 的所有 \(\frR\)-、\(\frS\)-自同態組成的環。作為 \(\frS\)-自同態環 \(\mA\) 的子環，\(\frB\) 具有一個在 \(\frS\) 中的反向同構表示。按照 §1, 2. 和 3.，\(\frB\) 由 \(\mA\) 中所有同 \(\frR\) 交換相容、也就是滿足關係 (2a) 的自同態組成。因而，\(\frB\) 是所有這樣的自同態的總體：它們與 \(\overline{\frR}\) 中的自同態逐個可交換，其中 \(\overline{\frR}\) 表示 \(\frR\) 在 \(\mA\) 中的像。由於 \(\frB\) 和 \(\overline{\frR}\) 的表示所涉及的是反向同構，而不僅僅是同態，這便是所要證明的事實。

\noindent\textbf{說明。} \(\frR\)-、\(\frS\)-自同態指的是一個對應 \(m_i\mapsto m_i'\)，使得藉助 \(m_i'\) 產生同一個表示 \(\frR^*\)；這些 \(m_i'\) 不一定構成 \(\frM\) 的一個 \(\frS\)-基，這同 \(\frB^*\) 中的矩陣不一定可逆這一事實相對應。當和 \(\sum m_i'\frS\) 不再是直和時，“表示”應按引申意義理解：
\[
        (m_1',m_2',\ldots,m_n')a=(m_1',\ldots,m_n')A
\]
仍然成立，但 \(A\) 不再由 \(a\) 唯一確定。

關於過渡定理還有如下說明：對角矩陣 \(E\cdot s\) 給出 \(\frS\) 的一個\emph{直接同構}表示，即 \(\frS\) 的“恆等”表示。這裡之所以是\emph{直接}同構，是因為 \(\frM\) 的 \(\frS\)-自同態環含有一個與 \(\frS\) 反向同構的子環，而這裡所涉及的是該子環的反向表示。事實上，\(\frM\) 的每個單生成元直和項 \(m_i\frS\) 作為右模都與 \(\frS\) 算子同構，所以它的 \(\frS\)-自同態環也與 \(\frS\) 反向同構（§1, 3., 推論 1）。

因此，由 \(\frR\) 與 \(\frS\) 唯一確定的從 \(\frT\) 到 \(\frS\) 中矩陣的對應，\emph{並不是環 \(\frT\) 的一個表示，而是把 \(\frR\) 的反向表示擴張為一個關於 \(\frS\) 算子同態的表示：} \(\sum r_i s_i\longmapsto \sum R_i s_i.\)
特別地，\(\frR\) 的反向表示也就\emph{關於 \(\frR\) 與 \(\frS\) 的交 \([\frR,\frS]\)——它位於 \(\frR\) 的中心內——具有算子同態性。}
\begin{center}
\textbf{§ 3.}\\[0.4em]
\textbf{關於除環的模。}
\end{center}

以下只考察在一般非交換除環中的表示；因此先彙集若干關於除環上線性形式模的簡單定理。

\noindent\textbf{1. 子模相對於全模給定基的正規基。} -- 設 \(A\) 為除環，\(\frN=x_1A+\cdots+x_nA\) 是秩為 \(n\) 的線性形式模；再設 \(\frL=z_1A+\cdots+z_\ell A\) 是秩為 \(\ell\le n\) 的子模。\emph{若 \(z_i\) 具有如下形式，則稱其為相對於 \(x\) 的正規基：}
\[
        z_i=x_i-(x_{\ell+1}\alpha_{i,\ell+1}+\cdots+x_n\alpha_{i,n}),
        \qquad (i=1,\ldots,\ell).
\]
\emph{每個模 \(\frL\) 在適當編號 \(x\) 後都有一個正規基。}

事實上，經適當編號後有
\[
        \frN=\frL+x_{\ell+1}A+\cdots+x_nA,
\]
從而
\[
        x_i\equiv x_{\ell+1}\alpha_{i,\ell+1}+\cdots+x_n\alpha_{i,n}
        \quad(\frL)\qquad (i=1,\ldots,\ell). \tag{2}
\]
因此，這 \(\ell\) 個元素 \(z_i=x_i-(x_{\ell+1}\alpha_{i,\ell+1}+\cdots)\) 位於 \(\frL\) 中，並且窮盡 \(\frL\)，因為它們與 \(x_{\ell+1},\ldots,x_n\) 一起構成 \(\frN\) 的一個基。

\noindent\textbf{2. 擴張模。} 仍設 \(A\) 為除環，\(P\) 為子除環。若秩為 \(n\) 的 \(A\)-模 \(\frN\) 含有一個同秩子模 \(\frM\)，則稱 \(\frN\) 為 \(P\)-模 \(\frM\) 的\emph{擴張模}；記作 \(\frN=\frM_A\)。若 \(\frM=y_1P+\cdots+y_nP\)，則有 \(\frN=\frM_A=y_1A+\cdots+y_nA.\) 反過來，對於每個 \(P\)-模 \(\frM=y_1P+\cdots+y_nP\)，其擴張模 \(\frM_A\) 在模同構意義下唯一存在。因為模 \(\overline y_1A+\cdots+\overline y_nA\) 含有一個與 \(\frM\) 同構的子模 \(\overline{\frM}\)，可將後者與 \(\frM\) 等同起來。

\noindent\textbf{推論 a)。} 若 \(z_1,\ldots,z_t\) 是 \(\frM\) 中關於 \(P\) 線性無關的元素，則它們在 \(\frN=\frM_A\) 中也關於 \(A\) 線性無關（因為可以把它們補成 \(\frM\) 的一個基，從而也補成 \(\frN\) 的一個基）。因此，\(\frM\) 中的每個 \(P\)-模 \(\frT=z_1P+\cdots+z_tP\) 都在 \(\frM_A\) 中生成一個擴張模 \(\frT_A=z_1A+\cdots+z_tA\)。

\noindent\textbf{推論 b)。} \(\frM\) 中的每個 \(P\)-模 \(\frT=z_1P+\cdots+z_tP\) 都是限制模，即其擴張模 \(\frT_A\) 與 \(\frM\) 的交；\(\frT=\frT_A\cap\frM\)。事實上，\(\frT\subseteq \frT_A\cap\frM\)；反過來，若 \(a\) 位於 \(\frT_A\cap\frM\)，則有 \(a=\sum z_i\alpha_i\)，於是 \(\frM\) 中的元素 \(a,z_1,\ldots,z_t\) 關於 \(A\) 線性相關，因而［推論 a)］也關於 \(P\) 線性相關；所以 \(a\) 位於 \(\frT\) 中。

\noindent\textbf{3. 關於不變模的定理。} 仍設 \(\frN=\frM_A\) 為秩為 \(n\) 的 \(P\)-模 \(\frM\) 的擴張模；再設對於 \(A\) 的一個環自同構群 \(\frG\)，\(P\) 是全不變除環。通過如下規定，把群 \(\frG\) 定義為 \(\frN=\frM_A\) 的算子域：
\[
        G\cdot m=m\quad\text{其中 }m\text{ 來自 }\frM\text{，而 }G\text{ 來自 }\frG,
\]
以及
\[
        G\cdot\sum m_i\alpha_i=\sum m_i\,G(\alpha_i)
        \quad\text{其中 }m\text{ 來自 }\frM\text{，而 }\alpha\text{ 來自 }A.
\]

\noindent\textbf{引理。} \emph{在這些規定下，\(\frM\) 由 \(\frN\) 中在 \(\frG\) 作用下保持不動的全部元素組成。} 事實上，若 \(x_1,\ldots,x_n\) 是 \(\frM\) 的一個 \(P\)-基，因而 \(w=x_1\alpha_1+\cdots+x_n\alpha_n\)，則對每個來自 \(\frG\) 的 \(G\)，由 \(G(w)=w\) 可推出也有 \(G(\alpha_i)=\alpha_i\)，從而根據假設，\(\alpha_i\) 位於 \(P\) 中。

\noindent\textbf{定理。} \emph{\(\frM\) 的子模 \(\frL\) 的擴張模 \(\frL_A\)，並且只有這些擴張模，是相對於作為算子域的 \(\frG\) 的容許子模。}

前一半是清楚的。反過來，設 \(\frL\) 相對於 \(\frG\) 是容許的，並設 \(z_1,\ldots,z_\ell\) 是 \(\frL\) 相對於 \(x\) 的一個正規基（見 1.）；這裡把 \(x\) 選作 \(\frM\) 的一個 \(P\)-基，因而 \(\frG\) 逐元素固定它。根據假設，
\[
        G(z_i)=x_i-\bigl(x_{\ell+1}G(\alpha_{i,\ell+1})
              +\cdots+x_nG(\alpha_{i,n})\bigr)
\]
對於每個來自 \(\frG\) 的 \(G\) 都是 \(\frL\) 的元素；所以，當 \(G\) 固定時，它可以由這些 \(z\) 線性表示。因此，由比較 \(x_1,\ldots,x_\ell\) 的係數可知：對每個來自 \(\frG\) 的 \(G\)，都有 \(G(z_i)=z_i\)。於是，根據引理，這些 \(z\) 位於 \(\frM\) 中，而 \(\frL\) 成為 \(\frM\) 中 \(P\)-模 \(\frT=z_1P+\cdots+z_\ell P\) 的\emph{擴張}，其中 \(\frT=\frL\cap\frM\)（根據 2）\footnote{藉助簡單的良序論證，本段的各定理即使在 \(\frM\) 關於 \(P\) 的秩變為無限時仍然成立（參見 G. Köthe：Ein Beitrag zur Theorie der kommutativen Ringe ohne Endlichkeitsvoraussetzung § 1, Gött. Nachr. 1931, S. 195--207）。}。

\begin{center}
\textbf{\S\ 4.}\\[0.5em]
\textbf{超復系統及其表示類。}
\end{center}

\noindent\textbf{1. 擴張定理。} 現在把 § 2 的表示定理同 § 3 的模定理結合起來。我們不再考察一般的環 \(\frR\)，而考察交換域 \(P\) 上帶單位元的\emph{超復系統}——\(S,T,\ldots\)——即
\[
        S=x_1P+\cdots+x_nP;
\]
並且不再考察任意表示環 \(\frS\)，而只考察除環 \(A,B,\ldots\)，若無另行說明，就考察中心為 \(P\) 的除環。在這種情形下，總存在一個乘積環，使 \(S\) 與 \(A\) 逐元素可交換，並且其交為 \(P\)；更確切地說，這裡涉及的是直積 \(S\times A\)，它同時成為 § 3 意義下 \(S\) 的擴張模 \(S_A\)
\[
        S_A=x_1A+\cdots+x_nA,
\]
因而也稱為擴張環。

在這種特化下，關於不變模的定理（§ 3, 3.）化為如下

\noindent\textbf{擴張定理：} \emph{\(S_A\) 中的每個雙側理想都是 \(S\) 中一個理想的擴張理想，更確切地說，是它同 \(S\) 的交的擴張。}

事實上，若所取群 \(\frG\) 是 \(A\) 的內自同構群，則作為 \(A\) 的中心，\(P\) 成為全不變域；而 \(A\) 與 \(S\) 逐元素可交換這一事實表明，\(\frG\) 按 § 3, 3. 擴張到 \(S_A\) 上，使得 \(S\) 成為全不變範圍。\(S_A\) 中的每個雙側理想 \(\mathfrak a\) 都是容許子群——因為 \(\alpha^{-1}\mathfrak a\alpha\leq\mathfrak a\)——所以根據 § 3, 3.，它是其交模 \(\mathfrak a\cap S\) 的擴張，而該交模成為 \(S\) 中的理想。

\noindent\textbf{附加說明：} 若 \(S\) 是交換的，則 \(S\) 成為 \(S_A\) 的中心。因為 \(S\) 位於中心中，並且相對於由 \(A\) 生成的內自同構，它是全不變範圍。\emph{更一般地，\(S\) 的中心也成為 \(S_A\) 的中心。}

\noindent\textbf{記號：} \(A_r,B_n,\ldots\) 以下始終表示 \(A,B,\ldots\) 上相應階數的矩陣環，即
\[
        A_r=\sum A c_{ik}=\sum c_{ik}A.
\]

下面將本質地使用一個推論，即簡單系統的擴張定理：若 \(S\) 是簡單系統\footnote{簡單系統照通常含義，是指\glqq{}雙側簡單\grqq{}。}，則 \(S_A\) 也簡單：\(S_A=\sum B c_{ik}=\sum c_{ik}B=B_t\)，其中 \(B\) 是對應的除環，其中心包含 \(P\)。這裡 \(A\)，因而 \(B\)，可以在 \(P\) 上具有無限秩；只假定 \(S\) 是超復的。若 \(S\) 與 \(A\) 有共同中心 \(P\)，則 \(P\) 也成為 \(S_A\) 的中心。

更一般地還成立：若 \(S\) 是簡單超復系統，則直積 \(S\times A_r\) 也簡單。因為 \(S\times A_r\) 與 \(S_A\times P_r\) 一致；若 \(S_A=B_t\)，它就等於 \(B_{tr}\)。

\noindent\textbf{2. 表示類。} 所謂超復系統的——反向或直接——表示，照通常理解，是指關於係數域 \(P\) 算子同態（§ 2 末尾）並且不同於零表示的表示。把過渡定理及其推論（§ 2, 3.）同 1. 中的擴張定理結合起來，得到如下

\noindent\textbf{關於表示類的定理：} \emph{設 \(S\) 是帶單位元的超復系統，以 \(P\) 為係數域，且 \(A\) 是中心為 \(P\) 的除環；那麼 \(S\) 在 \(A\) 中不同的不可約反向表示類的數目，等於 \(S_A\) 按其根基取商環後所得環中簡單右理想類的數目。——若 \(S\) 是簡單系統，則它在 \(A\) 中恰有一個不可約反向表示類和一個不可約直接表示類\footnote{\normalfont 參見《表示論》注 20），其中處理的是 \(A\) 為 \(S\) 的對應自同態除環的情形。}。}

事實上，根據過渡定理的推論（§ 2, 3.），簡單反向表示類與 \(S_A\) 上的簡單模類一一對應。由於 \(S_A\) 關於除環 \(A\) 具有有限秩，它滿足單側理想的極大條件和極小條件；於是斷言的第一部分直接由《表示論》§ 19 得出。若 \(S\) 簡單，則根據擴張定理，\(S_A\) 也簡單。由於極大條件和極小條件，它同時完全可約，因此只有一個簡單單側理想類；也就是說，\(S\) 在 \(A\) 中只有一個不可約反向表示類。與 \(S\) 反向同構的系統也隨 \(S\) 一起是簡單的，並且只有一個不可約反向表示類；所以 \(S\) 也只有一個簡單直接表示類。

\noindent\textbf{3. 秩關係：} 當 \(S\) 簡單時，\(S_A=\sum Bc_{ik}\) 無論關於 \(A\) 還是關於 \(B\) 都具有有限秩；這一事實給出秩關係：
\[
\begin{array}{rl}
(1)& n=rt\quad\text{其中}\quad n=(S:P),\quad t^2=(S_A:B),\quad r=\text{不可約反向表示的階}.
\end{array}
\]

因為 \(S_A\) 本身就是 \(S\) 在 \(A\) 中的反向表示模（該表示已經位於 \(P\) 中）。作為 \(B\) 上的矩陣環，\(S_A\) 分解為 \(t\) 個簡單右理想，即 \(S_A\)-模；這些模關於 \(A\) 的秩等於不可約反向表示的階 \(r\)。由於 \(n\) 同時也是 \(S_A\) 關於 \(A\) 的秩，秩關係隨之得出。

\paragraph{4. 當 \(A\) 關於 \(P\) 具有有限秩時，秩關係的加強。}
在這種情形下有三個關係：
\[
\begin{array}{rl}
(1)& (S:P)=n=rt,\\
(2)& (B:P)\cdot t=(A:P)\cdot r,\\
(3)& (S:P)\cdot(B:P)=(A:P)\cdot r^2.
\end{array}
\]
這裡，(2) 表示 \(S_A\) 中一個簡單右理想關於 \(P\) 的秩，按照 3.，分別用關於 \(B\) 或關於 \(A\) 的秩來表達；而 (3) 表示矩陣環 \(B_n\) 中一個簡單右理想的秩，並且由 (2) 藉助 (1) 乘以 \(r\) 得到。

\begin{center}
§ 5.

\textbf{表示定理在矩陣環上的應用。}
\end{center}

§ 4 中考察的 \(S\) 在 \(A\) 中的反向表示，也可以理解為 \(S\) 到 \(A_r\) 的一個子環中的直接同態；這裡 \(\overline A\) 表示與 \(A\) 反向同構的除環。\emph{以下各段的原則反過來是：把矩陣環 \(A_r\) 的研究歸結為反向同構除環 \(A\) 中的表示論。}

\noindent\textbf{1. 在 \(A_r\) 中的可約與不可約嵌入。} 設 \(A\) 與 \(\overline A\) 是兩個關於其中心 \(P\) 具有有限或無限秩的反向同構除環\footnote{一般說來，反向對應的除環或系統將以相應的希臘字母和拉丁字母表示。}。\emph{若 \(P\) 上的簡單系統 \(S\) 在 \(A\) 中具有一個 \(f\) 階不可約或可約反向表示，則稱 \(S\) 可以不可約方式或可約方式嵌入 \(A_f\)。} 若 \(S\) 可以不可約方式嵌入 \(A_r\)，則它可約地嵌入所有矩陣環 \(A_{rs}\)（\(s>1\)）中，因為它具有這些階數且僅具有這些階數的可約表示（§ 4, 2.）。因此，\(S\) 與 \(A_r\) 和 \(A_{rs}\) 的某些子環直接同構。該同構是 \(P\) 上恆等映射的延拓\footnote{如同表示的情形一樣，下文出現的同構始終是 \(P\) 上恆等映射的延拓——即按 \(P\) 的同構——以後不再逐次特別說明。}；根據 § 4, 3.，這裡 \(r\) 是 \(S\) 關於 \(P\) 的秩 \(n\) 的一個因子。

通過這種方式，可以得到各個 \(A_f\)（\(f=1,2,\ldots\)）中所有包含 \(P\)、並且關於 \(P\) 具有有限秩的簡單子環；因為每個這樣的子環都與 \(\overline A_f\) 的一個子環反向同構，從而容許在 \(A\) 中有一個可約或不可約反向表示。

\noindent\textbf{2. 關於內自同構的定理。} \emph{若 \(S^{(1)}\) 與 \(S^{(2)}\) 是 \(A_f\) 中兩個包含 \(P\)、並且關於 \(P\) 具有有限秩的簡單子環，且 \(S^{(1)}\) 與 \(S^{(2)}\) 按 \(P\) 同構，則這個同構由 \(A_f\) 的一個內自同構產生。}

因為 \(S^{(1)}\) 與 \(S^{(2)}\) 表示同一個簡單系統 \(S\) 在 \(A_f\) 中的兩種嵌入——一般是可約嵌入——從而表示 \(S\) 在 \(A\) 中的兩個 \(f\) 階直接表示；所以它們屬於 \(A\) 中同一個可約直接表示類；完成過渡的矩陣產生該自同構。

若 \(A\) 本身關於 \(P\) 具有有限秩，因而是超復的，則該定理化為熟知的命題：

\emph{\(A_f\) 中兩個包含中心 \(P\)、並且按 \(P\) 同構的簡單子系統，可通過 \(A_f\) 的一個內自同構相互變換。特別地，\(A_f\) 的每個自同構都是內自同構\footnote{\normalfont 當 \(A\) 關於 \(P\) 具有無限秩時，這一點不再成立（參見 Köthe 在注 5）所引的工作）。}。}

\noindent\textbf{3. 關於逐元素可交換子環的定理。}
這個定理仍然按照本段開頭提到的原則，由 § 2, 3. 中關於可交換矩陣的定理推出：

\emph{若 \(S\) 是 \(A_f\) 中一個包含 \(P\) 的簡單系統，\(R\) 表示 \(A_f\) 中與 \(S\) 逐元素可交換的全部元素，則 \(R\) 也成為矩陣環：\(R=\overline B_s\)。\(S_A\) 的對應除環 \(B\) 與 \(R\) 的對應除環 \(\overline B\) 反向同構。\(R\) 與 \(S\) 的交成為 \(S\) 的中心。當且僅當 \(S\) 不可約地嵌入 \(A_f\) 時，\(R\) 成為除環。}

事實上，根據 § 2, 3.，\(R\) 與 \(S\) 在 \(A\) 中的反向表示模 \(\frM\) 的自同態環直接同構；而這正是把 \(\frM\) 看作 \(S_A\)-模時的自同態環。若 \(\frM\) 分解為 \(s\) 個簡單 \(S_A\)-模，並且如 § 4, 1. 那樣令 \(S_A=\sum Bc_{ik}\)，則 \(R\) 同構於 \(\sum \overline Bc_{ik}\)，其中 \(B\) 與 \(\overline B\) 反向同構（因為根據 § 1, 3. 結尾，\(B\) 與 \(S_A\) 中簡單右理想、即簡單 \(S_A\)-模的自同態除環按 \(P\) 反向同構）。當且僅當 \(s\) 等於一，即 \(S\) 被不可約嵌入時，\(R\) 成為除環——\(R=\overline B\)。\(R\) 與 \(S\) 的交就是 \(S\) 的中心，這一點直接由 \(R\) 的定義得出。

\noindent\textbf{4. 超復矩陣環情形下的加強交換定理。} 在這種情形中，§ 4, 4. 的秩關係給出如下加強：

\emph{\(A_f\) 中包含 \(P\) 的簡單子系統——這裡 \(A\) 關於其中心 \(P\) 具有有限秩——組成若干對 \(S,\overline S\)，使得 \(\overline S\) 由與 \(S\) 逐元素可交換的全部元素組成，反之亦然。\(S_A\) 與 \(\overline S\) 的對應除環反向同構，\(S\) 與 \(\overline S_A\) 的對應除環也反向同構。\(S\) 與 \(\overline S\) 的交等於共同中心。當且僅當一個子環是除環時，另一個子環被不可約嵌入。\(S\) 與 \(\overline S\) 關於 \(P\) 的秩數之積等於 \(A_f\) 的秩。若 \(f=rs=\overline r\,\overline s\)——其中 \(S\) 可以不可約方式嵌入 \(A_r\)，而 \(\overline S\) 可以不可約方式嵌入 \(A_{\overline r}\)——則 \(S\) 與 \(\overline S\) 的對應除環同構於 \(A_g\) 中一對可交換子環，其中 \(f=g s\overline s\)。}

本質上只須證明關於秩數之積的斷言。先設 \(S\) 被不可約嵌入，因而 \(\overline S\) 是除環，並且與 \(B\) 反向同構——其中 \(S_A=B_t\)——則 § 4, 4. 的秩關係 (3) 給出乘積斷言，因為現在 \(f=r\)，而 \(\overline S\) 的秩等於 \(B\) 的秩。若 \(S\) 被可約嵌入，因而 \(\overline S=B_s\)，則 \(f=rs\)，把 (3) 乘以 \(s^2\) 即得該斷言。由此立即推出，\(S\) 由與 \(\overline S\) 逐元素可交換的全部元素組成，從而除最後一個斷言外，其餘全部都由 3. 得出。最後一個斷言則由兩次過渡到不可約嵌入得出。在 \(A_r\) 中——其中 \(f=rs\)——\(S\) 被不可約嵌入，\(S\) 與 \(R\) 成為成對可交換系統，而 \(R\) 同構於除環 \(B\)。另一方面，由於 \(S\) 與 \(\overline S\) 的反向對應關係，\(S\) 等於 \(\overline B_s\)——其中 \(\overline B\) 對於 \(\overline S\) 的意義，正如 \(B\) 對於 \(S\) 的意義——所以 \(R\) 可約地嵌入 \(A_r\)，並且不可約地嵌入 \(A_g\)，其中 \(r=g\overline s\)。在這裡，這對逐元素可交換的系統便分別同構於 \(B\) 與 \(\overline B\)。

\begin{center}
§ 6.

\textbf{簡單系統的伽羅瓦理論。}
\end{center}

§ 5, 4. 的加強交換定理，表達了以中心作為基礎域時簡單系統的伽羅瓦理論。我們先考察除環，然後考察一般簡單系統。

\noindent\textbf{1. 關於中心具有有限秩的除環的伽羅瓦理論。} \(A\) 的\emph{伽羅瓦群} \(\Gcal\) 定義為所有延拓中心 \(P\) 上恆等映射的自同構所成的群，即（§ 5, 2.）所有內自同構所成的群。因此，\(\Gcal\) 同構於 \(A^*/P^*\)，其中 \(A^*\) 和 \(P^*\) 分別表示 \(A\) 和 \(P\) 中非零元素的乘法群。因此，\(\Gcal\) 的子群 \(\mathfrak H\) 與 \(A^*\) 中包含 \(P^*\) 的子群 \(H^*\)，通過 \(\mathfrak H\sim H^*/P^*\) 一一對應。
若 \(H^*\) 添入零以後成為一個除環 \(H\)，則稱 \(\Gcal\) 的子群 \(\mathfrak H\) 是閉的。\(C\) 容許群 \(\mathfrak H\) 這一事實，等價於 \(C\) 與 \(H\) 逐元素可交換。這樣，交換定理（§ 5, 4.）化為如下

\begin{center}
\textbf{\srcspaced{非交換除環伽羅瓦理論的}\
\srcspaced{主定理：}}
\end{center}

\emph{位於 \(P\) 與 \(A\) 之間的除環 \(C\)，同 \(\Gcal\) 的閉子群 \(\mathfrak H\)，可以這樣一一對應：\(C\) 成為 \(\mathfrak H\) 的全不變除環，而 \(\mathfrak H\) 成為 \(C\) 的全不變群。}

\noindent\textbf{\srcspaced{2. 簡單系統的伽羅瓦理論。}}
若 \(A_f\) 是中心為 \(P\) 的簡單系統，則伽羅瓦群 \(\Gcal\) 仍是內自同構群，因而同構於
\[
        A_f^*/P^*,
\]
其中 \(A_f^*\) 現在表示 \(A_f\) 中正則元素的乘法群。若由 \(H^*\) 生成的子環 \(H\) 是簡單系統，並且 \(H^*\) 由 \(H\) 中的全部正則元素組成，則稱 \(\Gcal\) 的子群 \(\mathfrak H\sim H^*/P^*\) 是簡單閉的。這裡從交換定理到伽羅瓦理論的過渡由如下引理給出：

\noindent\textbf{\srcspaced{引理：}} \emph{\(A_f\) 中每個包含 \(P\) 的簡單系統，都由其正則元素所成的群 \(H^*\) 生成。} 若 \(P\) 含有無限多個元素，這是清楚的；因為用不定元形成的\glqq{}一般元素\grqq{}是正則的，並且通過適當特化，可以構成關於 \(P\) 的正則基元素。不過根據 Shoda，該引理一般成立\footnote{Shoda，注 3）所引的工作。那裡避免了不定元的引入。}。

藉助該引理，\(S\) 與一個簡單系統 \(\overline S\) 可交換，又等價於 \(S\) 容許由 \(\overline S\) 的正則元素誘導的自同構。因此，交換定理（§ 5, 4.）在這裡也化為如下

\noindent\textbf{\srcspaced{簡單系統伽羅瓦理論的主定理：}} \emph{\(A_f\) 中包含 \(P\) 的簡單系統 \(S\)，同 \(\Gcal\) 的簡單閉子群 \(\mathfrak H\)，可以這樣一一對應：\(S\) 成為 \(\mathfrak H\) 的全不變範圍，而 \(\mathfrak H\) 成為 \(S\) 的全不變群。}

\noindent\textbf{\srcspaced{3. 藉助延拓原則的證明。}}
對於除環的情形，我還要給出主定理的第二個證明；這個證明基於同構、亦即一階表示的延拓原則，並且由於不使用秩數計數，需要較少的有限性假設。特別地，由此可以看出，向無限秩除環 \(S\) 轉移將可能沿着什麼方向進行。這個證明也不使用“只有一個不可約反向表示類”這一事實，因此在交換情形中同樣成立（§ 8, 2.）。我先列出若干引理。

\noindent\textbf{\srcspaced{假設：}} 設 \(P\) 上的簡單系統 \(S\) 容許在 \(A\) 中有一個一階反向表示，因而是除環；這裡 \(A\) 關於其中心 \(P\) 可以具有有限或無限秩。給定 \(S_A\) 分解為 \(n\) 個簡單且算子同構的右理想（§ 4, 3）：
\[
        S_A=r_1+\cdots+r_n
            =e_1S_A+\cdots+e_nS_A
            =e_1A+\cdots+e_nA ;
\]
於是，由 \(e_i\) 生成的反向表示定義為
\[
        e_i s=e_i\sigma .
\]

\noindent\textbf{\srcspaced{引理 1：}} \emph{由 \(e_1,\ldots,e_n\) 生成的這 \(n\) 個反向表示彼此不同，因而是同一個表示類中的不同表示。因此，\(S\) 至少擁有與其秩一樣多的不同表示。}

為作證明，我們按照 § 2 末尾，把來自 \(S\) 的表示擴張為 \(S_A\) 的對應，使之關於 \(A\) 算子同態。這裡，這些對應由
\[
        e_i w=e_i\omega
\]
定義，其中對於 \(S_A\) 中每個 \(w\)，\(\omega\) 位於 \(A\) 中。若由 \(e_i\) 與 \(e_j\)，\(i\ne j\)，生成同一個表示，則對 \(S_A\) 中每個 \(w\) 都會有 \(e_iw=e_jw\)。但由於 \(e_ie_i=e_i\) 而 \(e_je_i=0\)，這是不可能的。

\noindent\textbf{\srcspaced{引理 2：}} \emph{若 \(T\) 是 \(P\) 與 \(S\) 之間的除環，且 \(s\) 是 \(S\) 關於 \(T\) 的秩，則 \(T\) 到 \(A\) 中的每個反向同構至少容許 \(s\) 個不同延拓。}

設這個反向同構、亦即 \(T\) 在 \(A\) 中的反向表示，由如下分解給出：
\[
        T_A=E_1T_A+\cdots+E_hT_A
            =E_1A+\cdots+E_hA;\qquad E_it=E_i\tau,
\]
其中 \(E_i\) 是冪等元（這並不限制一般性，因為由一個基元素 \(E_i\alpha\) 生成的表示，也可由 \(\alpha^{-1}\cdot E_i\alpha\)、即一個冪等元生成）。用 \(S\) 作右乘，得到
\[
        S_A=E_1S_A+\cdots+E_hS_A .
\]
這裡各個 \(E_iS_A\) \emph{關於 \(A\) 的秩全都等於 \(s\)}；事實上，如藉助 \(S\) 與 \(A\) 的可交換性把 \(S\) 的一個 \(T\)-基代入所顯示的，\(E_iS_A\) 的秩至多為 \(s\)：
\[
        S=Tw_1+\cdots+Tw_s
\]
\[
        E_iS_A=E_iT_Aw_1+\cdots+E_iT_Aw_s
              =E_iw_1A+\cdots+E_iw_sA .
\]
由於總和 \(S_A\) 的秩為 \(n=hs\)，每個 \(E_iS_A\) 的秩恰為 \(s\)。因此有
\[
        E_iS_A=r_{i1}+\cdots+r_{is}
              =e_{i1}A+\cdots+e_{is}A,
\]
其中由 \(e_{i1},\ldots,e_{is}\) 生成的這 \(s\) 個表示，根據引理 1 全都不同。它們全都是由 \(E_i\) 生成的 \(T\) 的表示的延拓；因為由 \(E_it=E_i\tau\) 可得
\[
        (e_{i1}+\cdots+e_{is})t=(e_{i1}+\cdots+e_{is})\tau
\]
從而由於直和分解，有 \(e_{ij}t=e_{ij}\tau\)。

\paragraph{定義。}
由 \(T\) 誘導的、從 \(S\) 到 \(A\) 的反向同構 \(\mathfrak I\) 的類劃分 \(\mathfrak I_T\)，定義如下：\(\mathfrak I\) 中全部且僅有那些在 \(T\) 上誘導同一同構的同構，被看作等價的。

\noindent\textbf{\srcspaced{主定理：}} \emph{若 \(\mathfrak I_T\) 是由 \(T\) 誘導的類劃分，則 \(T\) 是相對於這個類劃分的極大子除環。}

事實上，若 \(L\) 是位於 \(T\) 與 \(S\) 之間的中間除環，並且關於 \(T\) 的次數為 \(l\)，則根據引理 2，\(\mathfrak I_T\) 中的每個類在 \(\mathfrak I_L\) 中至少分裂為 \(l\) 個類。從主定理的這個表述過渡到通常表述，是通過把同構集合 \(\mathfrak I\) 的各元素同任意一個固定元素的逆相乘，使之進入自同構群；這時 \(\mathfrak I_T\) 的一個類轉入 \(T\) 的不變群。閉子群是全不變群這一斷言也包含在其中；只須對 \(\mathfrak H\sim H^*/P^*\)，把除環 \(H\) 作為中間除環 \(T\) 代入即可。
\begin{center}
§ 7.

\textbf{\srcspaced{相似代數的類。分裂域。}}
\end{center}

從現在起，只討論相對於其中心 \(P\) 秩有限的除環 \(A,\overline A\) 以及矩陣環 \(A_r\)；也就是說，只討論 \(P\) 上的簡單正規代數，以下簡稱為 \emph{\(P\) 上的代數}，或簡稱代數；於是 \(A,\overline A\) 是\emph{除法代數}。一個\emph{相似代數類}——\(A_r\sim A_s\)——彙集所有\emph{具有同一對應代數 \(A\)} 的 \(A_r\)。其中同構的 \(A\) 視為相同。

\noindent\textbf{\srcspaced{1. 代數類群。}} 若 \(A_r,B_s\) 是 \(P\) 上的代數，則它們的乘積（關於 \(P\) 的直積）也是如此；因為由 § 3 末尾可得：
\[
        A_r\times B_s=C_t,
\]
其中 \(C\) 是中心為 \(P\) 且在同構意義下唯一確定的除法代數。將 \(A_r,B_s\) 分別乘以 \(P\) 上的矩陣環可見，這一乘法由各個類唯一確定，因而這些類構成一個對乘法封閉的交換系統。此外還有

\noindent\textbf{\srcspaced{R. Brauer 定理：}} \emph{相似代數類關於直積構成一個阿貝爾群。}

事實上，這個系統具有一個單位類，它由 \(P\) 上的矩陣環，即與 \(P\) 相似的代數組成。並且，每個類 \((A)\) 都有逆類 \((A)^{-1}\)，它由與 \(\overline A\) 相似的代數組成，其中 \(\overline A\) 與 \(A\) 互反同構。這由 § 5, 3. 的交換定理在特化 \(S=A\) 時推出。因為 \(A\) 可以不可約地嵌入自身；所以得到 \(B=P\) 和 \(A_A=P_t\)\footnote{關於代數類群的更精細定理使用因子系統理論；這裡不作討論。應當指出，到此為止尚未出現關於交換域的考察；例如，秩 \((A:P)\) 是平方數這一事實既未使用，也未證明。}。

\noindent\textbf{\srcspaced{2. 一個代數類的分裂域。}}\quad 分裂域理論仍完全建立在 § 5 開頭所述的原則之上；交換擴張域在 \(A\) 或 \(\overline A\) 中表示。先給出

\medskip
\noindent\textbf{\srcspaced{輔助定理：}} \emph{若 \(A\) 是中心為 \(P\) 的除環，\(Z\) 是 \(P\) 的有限交換擴張域，則 \(A_Z\) 是中心為 \(Z\) 的簡單系統。} 因為根據 § 4, 1.，這對與 \(A_Z\) 互反同構的系統 \(Z_A\) 成立。因此，對 \(P\) 上每個由與 \(A\) 相似的代數組成的類 \((A)\)，\(Z\) 都產生一個\emph{擴張類} \((A)_Z\)，它由 \(Z\) 上與 \(A_Z\) 相似的簡單代數組成。

一個擴張類所對應的除法代數 \(D\)，可立即由 § 5, 3. 的交換定理得出：

\noindent\textbf{\srcspaced{關於擴張類所對應除法代數的定理：}} \emph{設 \((A)_Z\) 是一個擴張類，\(Z\) 表示 \(Z\) 在 \(A_r\) 中的一個不可約嵌入（表示）。則 \((A)_Z=(D)\)，其中 \(D\) 由 \(A_r\) 中同 \(Z\) 的每個元素都可交換的全部元素組成。} 只須在 § 5, 3. 中以 \(Z\) 代替那裡的簡單系統 \(S\)，並注意 \(Z_A\) 與 \(A_Z\) 互反同構。若涉及 \(Z\) 的可約嵌入，則同 \(Z\) 的每個元素都可交換的全部元素給出 \(D\) 上的一個矩陣環。

眾所周知，\(P\) 的交換擴張域 \(Z\) 稱為類 \((A)\) 的\emph{分裂域}，如果擴張類 \((A)_Z\) 完全分裂，即等於 \(Z\) 上的單位類，也就是 \(Z\) 上全部矩陣環的類。

\noindent\textbf{\srcspaced{關於分裂域刻畫的定理。}}
\emph{\(P\) 的一個有限交換擴張域 \(Z\)，當且僅當它在 \(A_r\) 中的不可約嵌入給出 \(A_r\) 的一個極大交換子域時，是類 \((A)\) 的分裂域。}

事實上，\(Z\) 是 \((A)\) 的分裂域這一性質等價於 \((D)=(Z)\)。因此，根據關於所對應除法代數的定理，不可約嵌入的 \(Z\) 必須是極大交換子域。若這一條件成立，則必有 \(D=Z\)，因為把 \(D\) 中的任一 \(a\) 添接到 \(Z\) 上都會生成一個交換擴張域。

\noindent\textbf{\srcspaced{說明。}}\quad 1. 同時可得，一個\emph{極大交換子域} \(Z\) 在\emph{不可約}嵌入時生成一個\emph{極大交換子環}。因為同 \(Z\) 的每個元素都可交換的全部元素構成一個除環。

2. 該定理同時給出分裂域的\emph{存在性}；因為所對應的除法代數 \(A\) 必定具有極大交換子域。

\medskip
\noindent\textbf{\srcspaced{3. 秩關係、指數、無限交換擴張。}}\quad § 5, 4. 的秩斷言首先給出一個熟知事實：簡單代數相對於其中心的秩是平方數；因為若 \(Z\) 在 \(A\) 中極大交換，則由於 \(A\) 中每個嵌入都是不可約的，有 \((Z:P)\cdot(Z:P)=(A:P)\)，從而 \((A:P)=m^2;\ (Z:P)=m\)，其中 \(m\) 是\emph{舒爾指數}；它同時表示所有可嵌入除法代數 \(A\) 本身的分裂域的次數。進一步，對一般分裂域的次數 \(n\)，有 \(n=m\cdot r\)\footnote{一般說來，也存在各種次數的“極小”分裂域，即沒有真子域仍是分裂域的分裂域。參見 Brauer--Noether 注 2）所引的工作。}；這例如可由 \((A_r:P)=m^2\cdot r^2\) 得到，也可由 § 4, 3. 的秩關係得到；後一種推導是因為指數 \(m\) 也等於 \(A_Z\) 的簡單分量數，而 \(A_Z\) 成為 \(Z\) 上秩為 \(m^2\) 的矩陣環。更一般地，若 \(A_\Lambda\sim D\)，\(L\) 是 \(\Lambda\) 在 \(A_r\) 中的不可約嵌入，且 \(d^2\) 是除法代數 \(D\) 關於其中心 \(L\) 的秩 \((D:L)\)，則有：\(l\cdot d=m\cdot r\)；因為根據 § 5, 4. 以及 4. 中關於所對應除法代數的定理，有：\((A_r:P)=(L:P)\cdot(D:P)\)，即 \(m^2r^2=l^2d^2\)。

由分裂域的存在性，對 \(P\) 的無限交換擴張域 \(\Omega\)（代數的或超越的），可得如下事實：

\noindent\emph{擴張類 \((A)_\Omega\) 是中心為 \(\Omega\) 的簡單代數類。特別地，若 \(\Omega\) 是代數閉的，則 \(A_\Omega\) 成為次數為 \(m\) 的矩陣環。因此，指數等於“絕對分量數”。}

事實上，若 \(Z\) 是 \(A\) 的分裂域，\(\bar\Omega\) 是 \(\Omega\) 與 \(Z\) 的合成域，則 \(A_{\bar\Omega}\) 成為 \(\bar\Omega\) 上的矩陣環，因而沒有根基且中心為 \(\bar\Omega\)。於是 \(A_\Omega\) 也沒有根基，並且其中心是 \(\Omega\)（因為 \(A_\Omega\) 中一個不屬於 \(\Omega\) 的元素也會位於 \(A_{\bar\Omega}\) 的中心），所以是除環。因此，\(A_\Omega\) 是 \(\Omega\) 上的簡單代數；若 \(\Omega\) 代數閉，則必為 \(\Omega\) 上的矩陣環。

\medskip
\noindent\textbf{\srcspaced{4. 可分分裂域的存在性。}}\quad \emph{每個類 \((A)\) 都具有可分分裂域，甚至具有可以嵌入 \(A\) 本身的這種分裂域}\footnote{參見 G. Köthe, Über Schiefkörper mit Unterkörpern zweiter Art über dem Zentrum, Journ. f. Math. 166 (1932), S. 182--184。這裡這個簡單得多的證明來自 M. Zorn 的一個觀察。}。

設基礎域 \(P\) 的特徵為 \(p\)；指數 \(m=s\cdot p^f\)，其中 \(s\) 與 \(p\) 互素。若 \(Z\) 是 \(A\) 的一個次數為 \(m\) 的分裂域，且 \(Z_0\) 是其中所含的第一類擴張——即 \((Z:Z_0)=p^f\)——則 \(A_{Z_0}\sim D\) 的指數等於 \(p^f\)（根據 3.）。現在要證明存在 \(Z_0\) 的一個可分擴張域 \(Z_1\)，使 \(D_{Z_1}\) 的指數成為 \(p^f\) 的真因子。根據 3.，當 \(\Omega\) 代數閉時，\(D_\Omega\) 沒有根基；所以 \(D\) 的約化判別式不為零。因此，\(D\) 中至少有一個元素 \(d\) 的（約化）跡不為零；這裡 \(d\) 不可能已經在基礎域 \(Z_0\) 中，因為 \(Z_0\) 中每個元素 \(\alpha\) 的跡都等於 \(p^f\alpha\)，因而為零。所以 \(Z_1=Z_0(d)\) 是一個真擴張域，並且是可分擴張域；\(D_{Z_1}\) 的指數成為 \(p^f\) 的真因子。有限次重複這個過程，就得到類 \((A)\) 的一個可分分裂域，其次數 \(m\) 與 \(A\) 的指數一致。

\begin{center}
\textbf{§ 8.}\\[0.35em]
\textbf{分裂域、析出域與伽羅瓦理論\\
在交換情形中。}
\end{center}

為了從相對於其中心為簡單的系統的分裂域過渡到任意簡單系統的分裂域，還須發展中心的分裂理論，並另行加入一種與 § 6, 3. 平行的伽羅瓦理論。（本段出現的全部域和系統都是交換的。）

\noindent\textbf{\srcspaced{1. 交換簡單系統的分裂域與析出域。}}\quad 設交換系統 \(Z\) 在 \(P\) 上簡單，即為一個域；\(\Omega\) 是 \(P\) 的一個代數閉擴張域。則眾所周知（《表示論》§21）：

\noindent\emph{\(Z_\Omega\) 當且僅當 \(Z\) 在 \(P\) 上可分（第一類擴張）時，仍為沒有根基的系統。}

因為當且僅當如此時，它才有與其相對於 \(\Omega\) 的秩同樣多的一階同構映射到 \(\Omega\) 中，因而也有同樣多的不同表示，而這裡表示與表示類相重合。

\(Z_\Omega\) 中可能出現根基，因而未必發生分解為絕對簡單分量的直和分解；這個事實引出如下

\noindent\textbf{\srcspaced{定義：}}\quad \(P\) 的一個擴張域 \(\Lambda\) 稱為分裂域，如果 \(Z_\Lambda\) 分裂為一階合成因子；換言之，如果所有絕對不可約表示都已在 \(\Lambda\) 中。\(P\) 的一個擴張域 \(T\) 稱為析出域，如果 \(Z_T\) 至少析出一個一階合成因子；換言之，如果 \(Z\) 在 \(T\) 中至少有一個絕對不可約表示。

分裂域和析出域可由下列命題刻畫：

\noindent\textbf{\srcspaced{定理：}} \emph{一個域 \(T\) 當且僅當含有一個與 \(Z\) 同構的子域 \(Z\) 時，是析出域；一個域 \(\Lambda\) 當且僅當含有一個子域 \(\Gamma\)，且 \(\Gamma\) 與屬於 \(Z\) 的伽羅瓦域同構時，是分裂域。}

這直接由藉助絕對不可約表示給出的分裂域和析出域的定義得出。特別地，類似於非交換情形，可得如下

\noindent\textbf{\srcspaced{推論：}} \emph{一個域 \(Z\) 當且僅當與 \(Z\) 同構時，才是極小析出域——即沒有真子域是析出域——（也就是說，它在 \(Z\) 中的不可約交換嵌入給出 \(Z\) 的一個極大交換子域）。一個極小析出域，當且僅當 \(Z\) 在 \(P\) 上正規、即為伽羅瓦域時，同時也是分裂域。}

這正是把簡單代數稱為相對於其中心正規的原因。在固定的、\(P\) 上的代數閉域 \(\Omega\) 內，只有\emph{唯一一個}極小分裂域，即屬於 \(Z\) 的伽羅瓦域；這與非交換情形中通常有無窮多個彼此不同構的極小分裂域\footnote{R. Brauer--E. Noether，同前引。}相對照。其原因在於，交換情形中的直和分解是唯一的，而非交換情形中的唯一性只到算子同構為止；或者——這說的是同一件事——在於交換情形中只有有限多個不同的絕對不可約表示，而不是像非交換情形那樣有無窮多個不同的絕對不可約表示，儘管後者屬於同一個類（參見 2.）。

\medskip
\noindent\textbf{\srcspaced{2. 可分域 \(Z/P\) 情形中伽羅瓦理論的超復奠基。}}\quad 與 § 6, 3. 完全類似，這裡對任意在 \(P\) 上可分的 \(Z\) 都有同構理論；當 \(Z\) 是伽羅瓦域時，就可以過渡到通常的自同構群。

\medskip
\noindent\textbf{\srcspaced{前提：}}\quad 設簡單系統 \(Z\) 是 \(P\) 的有限可分擴張域；\(\Omega\) 表示一個有限或無限的 \(P\) 上分裂域；\(Z=Z^{(1)}\) 是一個與 \(Z\) 同構的子域；\(\Gamma\) 是相應的伽羅瓦域。於是根據 1.（\(Z\) 可分！），有直和分解
\[
        Z=r^{(1)}+\cdots+r^{(n)}
        =e^{(1)}Z_\Omega+\cdots+e^{(n)}Z_\Omega
        =e^{(1)}\Omega+\cdots+e^{(n)}\Omega .
\]
由 \(e^{(i)}\) 生成的表示 \(Z\to Z^{(i)}\)，其中 \(Z^{(i)}\leqq\Gamma\)，由下式定義：
\[
        e^{(i)}z=e^{(i)}\zeta^{(i)}.
\]

\noindent\textbf{\srcspaced{輔助定理 1：}} \emph{由 \(e^{(1)},\ldots,e^{(n)}\) 生成的這 \(n\) 個表示全都不同。} 證明與 § 6, 3. 相同，也可由 1. 得出；因為它們屬於不同的表示類。

\noindent\textbf{\srcspaced{輔助定理 2：}} \emph{若 \(T\) 是 \(P\) 與 \(Z\) 之間的中間域，\(s\) 是 \(Z\) 相對於 \(T\) 的秩，則 \(T\) 到 \(\Omega\) 中的每個同構至少、因而恰好容許 \(s\) 個延拓。}

證明與 § 6, 3. 相同。這裡把“至少”加強為“恰好”，來自 \(Z_\Omega\) 分解的\emph{唯一性}；根據這一唯一性，\(Z\) 在 \(\Omega\) 中不是至少，而是恰好有 \(n\) 個同構（另參見 1. 末尾的說明）。

現在如 § 6, 3. 那樣，定義由 \(T\) \emph{誘導的類劃分} \(\mathfrak J_T\)，它是 \(Z\) 的有限同構集合 \(\mathfrak J\) 的一個類劃分——\emph{\(\mathfrak J\) 中全部且僅有那些在 \(T\) 上誘導同一同構的同構，在 \(\mathfrak J_T\) 中被視為等價}——並證明

\noindent\textbf{\srcspaced{主定理：}} \emph{若 \(\mathfrak J_T\) 是由 \(T\) 誘導的類劃分，則 \(T\) 是相對於這個類劃分的極大子域。}

從這裡出發，當 \(Z\) 是伽羅瓦域時，正如 § 6, 3. 一樣，可以通過與某個同構的逆相複合，過渡到自同構群和通常的表述。關於子群的相應斷言可由冪等元的考察得出，而這個考察本身也很有意義。

\medskip
\noindent\textbf{\srcspaced{3. \(Z_\Omega\) 的冪等元。}}\quad 令 \(S\) 遍歷把 \(Z\) 帶到各個 \(Z^{(i)}\) 中的 \(n\) 個代換，即同構 \(Z\to Z\) 與 \(Z\to Z^{(i)}\) 的複合，其中前者表示與 \(Z\to Z\) 互逆的同構。不必假定 \(Z\) 是伽羅瓦域。對 \(Z_Z\) 中的元素 \(a\)，把代換 \(S\) 定義為算子，規定它們在 \(Z\) 上誘導恆等映射；於是對每個 \(a\)，都定義了位於 \(Z_{Z^{(i)}}\) 中的共軛元 \(a^S\)。在這些規定下，有

\noindent\textbf{\srcspaced{定理：}} \emph{\(Z_\Omega\) 的 \(n\) 個冪等元 \(e^{(i)}\) 互為共軛：\(e^{(i)}=e^S\)，其中 \(e=e^{(1)}\)；它們位於 \(n\) 個共軛的擴張系統 \(Z_{Z^{(i)}}\) 中。}

事實上，應用 \(S\) 會把定義關係 \(ez=e\zeta^{(1)}\)、\(e^2=e\) 分別變為 \(e^Sz=e^S\zeta^{(i)}\)、\((e^S)^2=e^S\)。由於分解的唯一性，冪等元也唯一確定；所以 \(e^{(i)}=e^S\)。又因為 \(Z\) 是析出域，因而已經實現表示 \(ez=e\zeta\)，所以 \(e\) 已經位於 \(Z_Z\) 中，從而 \(e^{(i)}\) 位於 \(Z_{Z^{(i)}}\) 中。

若特別地 \(Z\) 是伽羅瓦域，則伽羅瓦理論主定理中\emph{關於子群的部分}立即推出；因為若 \(\mathfrak H\) 是伽羅瓦群 \(\Gcal\) 的子群，則由於直和的各個分量唯一確定，元素 \(\sum e^H\)（其中 \(H\) 遍歷 \(\mathfrak H\)）容許來自 \(\mathfrak H\) 的全部代換，而不容許任何除此以外的代換。由於群 \(\Gcal\) 是為 \(Z\) 定義的，這裡討論的是 \(Z\) 的伽羅瓦理論；藉助同構，\(Z\) 的相應理論也同時得到處理。

\paragraph{4. 冪等元同互補基的聯繫。}
這裡仍不必假定 \(Z\) 是伽羅瓦域。

\noindent\textbf{\srcspaced{定理：}} \emph{若 \(a_1,\ldots,a_n\) 是 \(Z\) 在 \(P\) 上的一組基，並置冪等元 \(e=a_1\beta_1+\cdots+a_n\beta_n\)——從而 \(e^S=a_1\beta_1^S+\cdots+a_n\beta_n^S\)——則 \(\beta_1^S,\ldots,\beta_n^S\) 表示互補基 \(b_1,\ldots,b_n\)（相對於 \(a_1,\ldots,a_n\)）在由 \(e^S\) 實現的映射下的像。}

事實上，在由 \(e^S\) 實現的映射 \(e^Sz=e^S\zeta^S\)，特別是 \(e^Sa_i=e^S\alpha_i^S\) 中，由於 \(e^S\cdot e^S=e^S\)、\(e^S\cdot e^R=0\)（\(S\ne R\)），有對應 \(e^S\mapsto1\)、\(e^R\mapsto0\)。所以得到：
\[
        1=\alpha_1^S\beta_1^S+\cdots+\alpha_n^S\beta_n^S;\qquad
        0=\alpha_1^S\beta_1^R+\cdots+\alpha_n^S\beta_n^R\quad(S\ne R),
\]
或者寫成矩陣，即互補基的定義方程
\[
\begin{pmatrix}
\alpha_1 & \cdots & \alpha_n\\
\cdots & \cdots & \cdots\\
\alpha_1^S & \cdots & \alpha_n^S\\
\cdots & \cdots & \cdots
\end{pmatrix}
\cdot
\begin{pmatrix}
\beta_1 & \cdots & \beta_1^R & \cdots\\
\cdots & \cdots & \cdots & \cdots\\
\beta_n & \cdots & \beta_n^R & \cdots
\end{pmatrix}
=
E.
\]
過渡到跡定義，是因為對 \(a_i=\sum_S e^S\alpha_i^S\) 和 \(b_i=\sum_S e^S\beta_i^S\)，該矩陣關係在交換矩陣次序後變為
\[
        \Sp(a_i b_i)=1,\qquad \Sp(a_i b_k)=0\quad\text{當 }i\ne k\text{ 時\srcfnmark{21)}}.
\]
\srcfntext{21)}{參見《表示論》§25。}
其中 \(a_i\) 與 \(b_i\) 位於 \(Z\) 中，因為它們容許全部代換 \(S\)（輔助定理 §3, 3.）。\srcfn{22)}{互補基與冪等元之間的關係最早見於 Dedekind, Zur Theorie der aus \(n\) Haupteinheiten gebildeten komplexen Größen；參見《全集》第二卷，第 4--5 頁。}

互補基的逆變性，在這裡來自如下事實：冪等元 \(e^S\) 是 \(Z\) 的不變量，並且與所取的基無關；也就是說，當過渡到 \(Z/P\) 的另一組基 \(\bar a_1,\ldots,\bar a_n\) 時，所有 \(a_1\beta_1^S+\cdots+a_n\beta_n^S\) 都轉入自身。

\begin{center}
\textbf{§ 9.}\\[0.8ex]
\textbf{任意系統的分裂域與析出域。}
\end{center}

\noindent\textbf{\srcspaced{1. 定義以及向簡單情形的歸約。}}
在一般情形中，與 § 8 所給定義相應的有如下

\noindent\textbf{\srcspaced{定義：}} 設 \(S\) 是 \(P\) 上的超復系統。\(P\) 的一個（交換）擴張域 \(\Lambda\) 稱為 \(S\) 的分裂域，如果 \(S_\Lambda\) 在一條由單側理想構成的合成列中分裂為絕對簡單合成因子；換言之，如果 \(S\) 在 \(\Lambda\) 中的全部不可約表示已經都是絕對不可約的。\(P\) 的一個擴張域 \(T\) 稱為析出域，如果 \(S_T\) 至少析出一個絕對簡單合成因子，也就是說，如果至少有一個在 \(T\) 中不可約的表示已經是絕對不可約的。

這些定義顯然同時包含 § 8 中對交換情形給出的定義和 § 7 中對簡單代數給出的定義；後一點是因為這裡 \(A_\Lambda\) 在中心的任一交換擴張下都是完全可約的，因而合成因子與直和分解的分量一致。但是，中心上的全矩陣環，而且只有這些環，才給出絕對簡單分量，因而也給出絕對不可約表示。又因為 \(A_\Lambda\) 是雙側簡單的，即分解為算子同構的分量，所以析出一個絕對簡單分量就已經意味着完全分裂：析出域與分裂域重合。

由這些定義立即得到如下事實：\emph{\(S\) 的分裂域，是為每個在 \(P\) 中不可約的表示各取一個分裂域後所得的合成域；\(S\) 的析出域，是任一在 \(P\) 中不可約表示的析出域。} 由於簡單系統——關於根基的商環的分量——與在 \(P\) 中不可約的表示之間有一一對應，所以只須\emph{考察簡單系統}。這裡的結果將由 § 7 與 § 8 組合得到。


\noindent\textbf{\srcspaced{2. 簡單系統的分裂域與析出域。}} 我們先考察簡單系統 \(A\) 的中心 \(Z\) 在 \(P\) 上\emph{可分}的情形。由 § 8, 1. 與 § 8, 3. 可得：

\emph{與 \(Z_\Omega\) 的直和分解相應的 \(A_\Omega\) 的雙側分解——其中 \(\Omega\) 表示 \(Z\) 的一個分裂域——給出 \(A\) 的 \(n\) 個共軛且與 \(A\) 同構的分量 \(e^{(i)}A\)，它們是各自中心 \(e^{(i)}Z^{(i)}=e^{(i)}Z\) 上的簡單代數。這裡，把代換 \(S\) 定義為 \(A_Z\) 的算子，規定它們在 \(A\) 上誘導恆等映射。}（參見 § 8, 3.）

事實上，根據 § 8, 3.，冪等元互為共軛；所以按照對 \(A\) 的代換定義，分量 \(e^{(i)}A\) 也同樣共軛。由於 \(A\) 是雙側簡單的，\(e^{(i)}A\) 與 \(A\) 環同構。又因為 \(e^{(i)}Z=e^{(i)}Z^{(i)}=e^{(i)}Z_{Z^{(i)}}\)，有 \(e^{(i)}A=e^{(i)}A_{Z^{(i)}}\)；所以中心與係數域重合，這些分量是正規代數。

§ 7, 2. 的結果進一步給出：

\emph{若 \(A\) 是 \(P\) 上的簡單系統，且其中心 \(Z\) 在 \(P\) 上可分，則 \(A_\Omega\) 也沒有根基，因而完全可約；這裡 \(\Omega\) 可以表示任一有限或無限擴張域。}

因為根據 § 7, 2.，\(e^{(i)}A_{Z^{(i)}}\) 在每次係數擴張下仍是簡單的，因而沒有根基。所以直和也同樣如此；而這個直和就是 \(A_\Omega\)，或者——當 \(\Omega\) 不是 \(Z\) 的分裂域時——由 \(A_\Omega\) 再作係數擴張得到。進一步，由 § 7, 2. 可得

\begin{center}
\textbf{\srcspaced{析出域與分裂域的刻畫。}}
\end{center}

\emph{若 \(A\) 是中心 \(Z\) 可分的除法代數，則析出域的不可約嵌入，恰由 \(A_r\) 中全部且僅有的包含 \(Z\) 的極大交換子域給出；分裂域全部且僅是一個析出域同中心的一個分裂域的合成域，特別是同屬於該中心的伽羅瓦域的合成域。即使中心不是伽羅瓦域，一個極小析出域也可能是分裂域。}

事實上，一個析出域必須包含一個與 \(Z\) 同構的域；關於析出域的嵌入斷言，現由前述事實和 § 7, 2. 推出。一個分裂域還必須包含 \(Z\) 的一個分裂域 \(\Omega\)。但是，\(\Omega\) 上的每個析出域同時也是分裂域，因為分量 \(e^{(i)}A_\Omega\) 都是簡單的，並且具有一個與 \(\Omega\) 同構的中心。若 \(Z\) 是伽羅瓦域，則極小析出域與分裂域重合。但即使在非伽羅瓦情形中，也可能與交換情形不同，存在同時為分裂域的極小析出域；只要這樣一個極小析出域已經包含屬於 \(Z\) 的伽羅瓦域，就會出現這種情形［例：中心與 \(P(\sqrt[3]{2})\) 同構的四元數除環，其中 \(P\) 是有理數域；屬於 \(P(\sqrt[3]{2})\) 的伽羅瓦域就是一個極小析出域和分裂域］。

最後，若涉及\emph{不可分中心，則 \(Z_\Omega\)，從而 \(A_\Omega\)，成為具有根基的系統。} 設 \(\mathfrak C\) 是 \(Z_\Omega\) 的根基，\(\mathfrak c\) 是 \(A_\Omega\) 中的擴張理想。於是，\(Z_\Omega/\mathfrak C\) 有一個分解為共軛分量的直和分解，與 § 8, 2. 完全類似；這又如上面一樣給出 \(A_\Omega/\mathfrak c\) 的直和分解，使得各分量 \(\bar e^{(i)}A\) 與 \(A\) 同構，並以 \(\bar e^{(i)}Z^{(i)}\) 為中心。因此，關於析出域和分裂域的定理仍然成立。進一步，秩的計數表明，與 \(\mathfrak C\) 的一條合成列相應的合成因子，對於 \(A_\Omega/\mathfrak c\) 成為算子同構，而不僅是同態。

用不可約表示來表述，可得如下\textbf{\srcspaced{概括：}} \emph{若給定一個超復系統的、在 \(P\) 中不可約的表示，則這個表示當且僅當相應的中心在 \(P\) 上可分時，才成為絕對完全可約的。在任何情形中，該表示的分裂域與析出域都可通過如上嵌入相應的簡單系統來刻畫。特別地，一個析出域的最低次數等於中心的次數與相應代數類在該中心上的指數之積；每個析出域的次數都是這個數的倍數。該表示——在合成列的意義下——分裂為若干不同但共軛的絕對不可約表示類，其數目等於中心所含最大可分子域的次數。每個類出現的重數，等於指數與中心相對於這個可分擴張的次數之積。}

對於完美基礎域，由於只有可分擴張域，我們由此回到 I. Schur 的已知結果，並補上了一個已見於 R. Brauer 的嵌入刻畫。

\begin{center}
（1932 年 6 月 8 日收稿。）
\end{center}

% --- RA10 appended body from N41_42_DE.tex ---
\clearpage

\setcounter{footnote}{0}

\end{document}
